\section{Partitions and the Rogers-Ramanujan identities}\label{sec:section_4}

A {\it partition (of length $r$)} of a postive integer $n$ is a
non-decreasing sequence $\lambda=(\lambda_1,\ldots, \lambda_r)$ of
positive integers $\lambda_i$, $1\leq i \leq r$, such that $$\lambda_1+
\cdots + \lambda_r =n.$$ The integers $\lambda_i$ are called the {\it
  parts of the partition $\lambda$}. We will denote the number of partitions of
$n$ by $p(n)$, with $p(0):=1$. In the
following we collect a few facts about integer partitions which will
be used in the subsequent sections. For an introduction to this topic
we refer for example to \cite{wilf_lecturesonis}; an extensive
treatment can be found in \cite{andrews_partitions}.

\begin{Prop}\label{prop:generatingfunction_partitions}
  The generating series of the partition function $p$ has the
  following infinite product representation: $$\sum_{i=0}^\infty p(n)t^n = \prod_{i\geq
    1}\frac{1}{1-t^i}.$$
\end{Prop}

Note, that this is precisely the series $\Hh$ which we have obtained as the
Hilbert-Poincar\'{e} series of the graded algebra $k[y_1, y_2,
\ldots]$ where the grading is given by $\wt\, y_i =i$. More generally,
the arc Hilbert-Poincar\'{e} series of an $d$-dimensional variety at a smooth point was given by
$\Hh^d$. This leads us to expect a connection between the
Hilbert-Poincar\'{e} series of arc algebras and partitions.\\

The following result is known in the literature as the {\it (first)
 Rogers-Ramanujan identity}. For a classical proof and an account of
its history, see Chpt.\ 7 of \cite{andrews_partitions}.

\begin{Thm}[Rogers-Ramanujan identity]\label{thm:rogersramanujan}
  The number of partitions of $n$ into parts congruent to 1 or 4 modulo 5
  is equal to the number of partitions of $n$ into parts that are neither
  repeated nor consecutive.
\end{Thm}

Many proofs of this identity can be found in the literature. See for instance \cite{andrews_proofs} for an overview of some of them. The Rogers-Ramanujan identity was generalized by Gordon. The statement that we need is the following, see Theorem 7.5 from \cite{andrews_partitions}.

\begin{Thm}\label{thm:gordon}
Let $k\geq 2$. Let $B_k(n)$ denote the number of partitions of $n$ of the form $(\lambda_1,\ldots,\lambda_r)$, where $\lambda_j - \lambda_{j+k-1}\geq 2$ for all $j \in \{ 1, \ldots, r-k+1\}$. Let $A_k(n)$ denote the number of partitions of $n$ into parts which are not congruent to $0, k$ or $k+1$ modulo $2k+1$. Then $A_k(n) = B_k(n)$ for all $n$.
\end{Thm}

The analytic counterpart of Theorem~\ref{thm:rogersramanujan} can be formulated as (see Corollary 7.9 in \cite{andrews_partitions}):

\begin{Cor}[Rogers-Ramanujan identity, analytic
  form]\label{cor:rogersramanujan_analytic}
Theorem~\ref{thm:rogersramanujan} is equivalent to the identity
$$1 + \frac{t}{1-t} + \frac{t^4}{(1-t)(1-t^2)} +
\frac{t^9}{(1-t)(1-t^2)(1-t^3)} + \cdots = \prod_{\substack{i\geq 1\\i\equiv 1,4\bmod 5}}\frac{1}{(1-t^i)}.$$
\end{Cor}

The analytic analogue of Theorem~\ref{thm:gordon} is somewhat more involved and we will not formulate it here. The interested reader can find it in Andrews' book.